\documentclass[12pt, a4paper]{ctexart}
\usepackage{fancyhdr, amsmath, amssymb, geometry, enumitem}
\geometry{top=2.5cm, bottom=2.5cm, left=2cm, right=2cm}
\setCJKmainfont{Hiragino Sans GB}
\setCJKsansfont{Heiti SC}
\pagestyle{fancy}
\fancyhf{}
\fancyhead[L]{\small 培优专题·函数与动点综合}
\fancyhead[R]{\small 测试卷}
\fancyfoot[C]{\thepage}
\renewcommand{\headrulewidth}{0.4pt}

\begin{document}

\begin{center}
{\Large\bfseries 培优专题·函数与动点综合\quad 测试卷}\\[6pt]
{\normalsize 总分：150分\quad 考试时间：120分钟}
\end{center}
\vspace{6pt}
\noindent\rule{\textwidth}{0.5pt}

\section*{一、单项选择题（本大题共8小题，每小题5分，共40分。每小题只有一个正确选项。）}

\begin{enumerate}

\item 动点$P$在单位圆$x^2+y^2=1$上运动，定点$A(1,0)$，则$|PA|$的最大值为\ （\quad）
\begin{enumerate}[label=\Alph*., itemsep=0pt, topsep=2pt]
\item $1$ \quad \item $\sqrt{2}$ \quad \item $2$ \quad \item $\sqrt{3}$
\end{enumerate}

\item 已知动点$P$在抛物线$y^2=4x$上运动，$F$为焦点，$A(3,0)$为定点，则$|PF|+|PA|$的最小值为\ （\quad）
\begin{enumerate}[label=\Alph*., itemsep=0pt, topsep=2pt]
\item $3$ \quad \item $4$ \quad \item $5$ \quad \item $2\sqrt{3}$
\end{enumerate}

\item 函数$f(x)=|x-1|+|x-2|+|x-3|$的最小值为\ （\quad）
\begin{enumerate}[label=\Alph*., itemsep=0pt, topsep=2pt]
\item $1$ \quad \item $2$ \quad \item $3$ \quad \item $4$
\end{enumerate}

\item 动点$P$在椭圆$\dfrac{x^2}{4}+\dfrac{y^2}{3}=1$上，$F_1$、$F_2$为焦点，则$\overrightarrow{PF_1}\cdot\overrightarrow{PF_2}$的取值范围是\ （\quad）
\begin{enumerate}[label=\Alph*., itemsep=0pt, topsep=2pt]
\item $[-1,2]$ \quad \item $[2,3]$ \quad \item $[0,4]$ \quad \item $[1,3]$
\end{enumerate}

\item 函数$f(x)=\dfrac{x^2+2}{x^2+1}$（$x\in\mathbb{R}$）的值域为\ （\quad）
\begin{enumerate}[label=\Alph*., itemsep=0pt, topsep=2pt]
\item $[1,2)$ \quad \item $[1,2]$ \quad \item $(1,2]$ \quad \item $(1,2)$
\end{enumerate}

\item 动点$P$从原点出发，沿$x$轴正方向以速度$1$运动，同时动点$Q$从点$(0,1)$出发，沿$y$轴正方向以速度$2$运动，则$|PQ|$关于时间$t$（$t\geq 0$）的函数的最小值为\ （\quad）
\begin{enumerate}[label=\Alph*., itemsep=0pt, topsep=2pt]
\item $\dfrac{\sqrt{5}}{5}$ \quad \item $\dfrac{2\sqrt{5}}{5}$ \quad \item $\dfrac{1}{2}$ \quad \item $\dfrac{\sqrt{2}}{2}$
\end{enumerate}

\item 已知函数$f(x)=x^2-2ax+1$在区间$[1,+\infty)$上的最小值为$g(a)$，则$g(a)$的最大值为\ （\quad）
\begin{enumerate}[label=\Alph*., itemsep=0pt, topsep=2pt]
\item $0$ \quad \item $1$ \quad \item $-1$ \quad \item $2$
\end{enumerate}

\item 已知动点$P$在曲线$y=\sqrt{1-x^2}$（$0\leq x\leq 1$）上运动，定点$A(1,1)$，则$|PA|$的最小值为\ （\quad）
\begin{enumerate}[label=\Alph*., itemsep=0pt, topsep=2pt]
\item $\sqrt{2}-1$ \quad \item $\sqrt{3}-1$ \quad \item $1$ \quad \item $\sqrt{5}-1$
\end{enumerate}

\end{enumerate}

\section*{二、多项选择题（本大题共4小题，每小题5分，共20分。每小题有多个正确选项，全部选对得5分，部分选对得2分，有选错得0分。）}

\begin{enumerate}[resume]

\item 已知函数$f(x)=|x-a|+|x-b|$（$a<b$），下列结论正确的是\ （\quad）
\begin{enumerate}[label=\Alph*., itemsep=0pt, topsep=2pt]
\item $f(x)$的最小值为$b-a$
\item $f(x)$在$(-\infty,a]$上单调递减
\item $f(x)$在$[b,+\infty)$上单调递增
\item $f(x)$在$[a,b]$上为常数函数
\end{enumerate}

\item 动点$P$在双曲线$\dfrac{x^2}{4}-\dfrac{y^2}{3}=1$的右支上运动，$F_1$、$F_2$分别为左、右焦点，下列结论正确的是\ （\quad）
\begin{enumerate}[label=\Alph*., itemsep=0pt, topsep=2pt]
\item $|PF_1|-|PF_2|=4$
\item 当$P$为右顶点时，$|PF_2|$取最小值$\sqrt{7}-2$
\item $\overrightarrow{PF_1}\cdot\overrightarrow{PF_2}\geq 0$
\item $\triangle PF_1F_2$面积的最小值为$3\sqrt{3}$
\end{enumerate}

\item 关于函数$g(x)=x+\dfrac{4}{x}$（$x>0$），下列结论正确的是\ （\quad）
\begin{enumerate}[label=\Alph*., itemsep=0pt, topsep=2pt]
\item $g(x)$在$(0,2]$上单调递减
\item $g(x)$在$[2,+\infty)$上单调递增
\item $g(x)$的最小值为$4$
\item 方程$g(x)=k$（$k>4$）有两个不相等的正实数解
\end{enumerate}

\item 已知$P$为圆$(x-2)^2+y^2=1$上的动点，$A(0,0)$，$B(4,0)$，下列结论正确的是\ （\quad）
\begin{enumerate}[label=\Alph*., itemsep=0pt, topsep=2pt]
\item $|PA|$的最小值为$1$
\item $|PB|$的最大值为$3$
\item $\overrightarrow{PA}\cdot\overrightarrow{PB}$的最大值为$6$
\item $|PA|^2+|PB|^2$为定值
\end{enumerate}

\end{enumerate}

\section*{三、填空题（本大题共4小题，每小题5分，共20分。）}

\begin{enumerate}[resume]

\item 动点$P$在圆$x^2+y^2=4$上运动，定点$A(4,0)$，线段$PA$的中点为$M$，则$M$的轨迹方程为\underline{\qquad}。

\item 已知函数$f(x)=\dfrac{x^2+3}{x^2+1}$（$x\in\mathbb{R}$），则$f(x)$的最小值为\underline{\qquad}。

\item 动点$P$在椭圆$\dfrac{x^2}{9}+\dfrac{y^2}{4}=1$上，$A(0,2)$为上顶点，则$|PA|$的最大值为\underline{\qquad}。

\item 已知正数$x$、$y$满足$x+y=1$，则$\dfrac{1}{x}+\dfrac{4}{y}$的最小值为\underline{\qquad}。

\end{enumerate}

\section*{四、解答题（本大题共6小题，共70分。解答应写出文字说明、证明过程或演算步骤。）}

\begin{enumerate}[resume]

\item （本小题满分10分）\\
动点$P$在抛物线$y^2=4x$上运动，$A(a,0)$（$a>0$）为$x$轴上的定点。
\begin{enumerate}[(1)]
\item 当$a=2$时，求$|PA|$的最小值；
\item 当$|PA|$的最小值为$2$时，求$a$的值。
\end{enumerate}

\item （本小题满分10分）\\
已知函数$f(x)=|2x-1|+|x+2|$。
\begin{enumerate}[(1)]
\item 求$f(x)$的最小值及取最小值时$x$的值；
\item 解不等式$f(x)\leq 5$。
\end{enumerate}

\item （本小题满分12分）\\
在平面直角坐标系中，$A(-2,0)$，$B(2,0)$，动点$P$满足$|PA|+|PB|=6$。
\begin{enumerate}[(1)]
\item 求$P$的轨迹方程；
\item 求$\overrightarrow{PA}\cdot\overrightarrow{PB}$的取值范围；
\item 求$\triangle PAB$面积的最大值。
\end{enumerate}

\item （本小题满分12分）\\
已知函数$f(x)=x^2-2x+2$，$g(x)=x^2+ax+1$。
\begin{enumerate}[(1)]
\item 若$f(x)\geq g(x)$对一切$x\in\mathbb{R}$恒成立，求$a$的取值范围；
\item 若存在$x_0\in[0,3]$使得$f(x_0)=g(x_0)$，求$a$的取值范围；
\item 若对一切$x\in[0,3]$，都有$f(x)\geq g(x)$，求$a$的取值范围。
\end{enumerate}

\item （本小题满分13分）\\
已知圆$C$：$(x-1)^2+y^2=1$，动点$P$在圆$C$上运动，$A(4,0)$为圆外定点。
\begin{enumerate}[(1)]
\item 求$|PA|$的最大值和最小值；
\item 求$\dfrac{y_P}{x_P-4}$的取值范围（$P(x_P,y_P)$为圆上一点）；
\item 过$A$作圆$C$的切线，求切线方程。
\end{enumerate}

\item （本小题满分13分）\\
已知椭圆$\dfrac{x^2}{4}+y^2=1$，$F_1$、$F_2$分别为左、右焦点，$A(1,0)$为椭圆内一点。动弦$BC$过点$A$。
\begin{enumerate}[(1)]
\item 求$\triangle F_1BC$周长的最小值；
\item 求$\overrightarrow{BF_1}\cdot\overrightarrow{CF_2}$的取值范围；
\item 若$BC$为过$A$且垂直于$x$轴的弦，求$\triangle F_1BC$的面积。
\end{enumerate}

\end{enumerate}

\end{document}
